\section{Continuous investment: primitive proofs}\label{sec:proofs95}
\subsection{Curvature, the action constraint, and state moduli}
\begin{proof}[Proof of Proposition~\ref{prop:curvature95}]
For $x\in X$, $0\leq a\leq1/4$ and $|z|\leq1/32$, the transition remains in $X$. Indeed, the state part is increasing in both of its arguments, its maximum is $1/16+1/2+1/8$, and $b_ja\leq1/8$. Its minimum after the shock is at least $1/32$, and its maximum is at most $27/32$. Thus every derivative bound below is used on its stated domain.

Both $S$ and $g$ are convex and continuously differentiable; their squared hinges have positive semidefinite second derivative on each smooth region. Direct differentiation gives
\[
 |\partial_jg|\leq\frac{10}{d},\qquad
 |\partial_jS|\leq\frac{15}{2d},\qquad
 \nabla^2S\succeq\frac4d I.
\]
The first bound adds $5/d$ from the target term, $1/d$ from the cyclic difference term, and $4/d$ from the hinge. Only $-x_jx_{j+1}/16$ in $F_j$ has nonzero second derivative. For a joint direction $(v,w)$, the possibly negative part of the second derivative of $g(F)$ is
\[
 -\frac18\sum_j\partial_jg(F)v_jv_{j+1}
 \ \geq\ -\frac{5}{4d}\|v\|_2^2,
\]
since $\sum_j|v_jv_{j+1}|\leq\|v\|_2^2$. The remaining term is the quadratic form of the positive semidefinite Hessian of $g$ applied to the full state-action Jacobian. In particular, it includes the mixed state-action terms and is not discarded selectively. Since $h''(a)=2+48a^2\geq2$, every fixed-shock objective satisfies \eqref{eq:curvature95}. Integrating against a fixed feasible risk density preserves this bound. Taking the supremum over those densities preserves the same strong convexity. The bound extends across the hinge by continuity of the first derivative, or by integrating the almost-everywhere second-derivative inequality along a segment.

The graph $\{(x,a):x\in X,0\leq a\leq A(x)\}$ is convex because $A$ is affine. Partial minimization of the jointly convex objective over this graph proves convexity of $V_1$. Compactness and continuity supply an attained minimum. The positive capacity permits constraint qualification in the action variable.

For a fixed feasible density, the supporting strong-convexity inequality gives, for every $(y,u)$,
\[
 Q_1(y,u)\geq W+p_x^\top(y-x_0)+p_a(u-a_0)+(u-a_0)^2.
\]
For feasible $u$, adding $\lambda(u-A(y))-\nu u\leq0$ to the right side preserves a lower bound. Minimize this augmented quadratic over all real $u$, an enlargement of the feasible interval. Completing the square gives \eqref{eq:affine95}. This proof uses neither equality in the risk supremum nor exact stationarity. Interval implementations bound the complete expression outwards, including its residual square and its slope.

For completeness we bound true value subgradients, rather than the potentially large slopes of arbitrary poor lower witnesses. A transition Jacobian column has nonnegative entries summing to at most $11/16$. Therefore any state subgradient of a risk-maximizing one-date objective has coordinate magnitude at most
\[
 \frac1d\left(\frac{15}{2}+\beta\frac{55}{8}\right).
\]
At an upper-bound action, the multiplier can be chosen with
$0\leq\lambda\leq\beta(10/d)\sum_j b_j=15\beta/4$:
the nonnegative derivative of $h$ cannot increase the negative derivative which this multiplier offsets. At a lower-bound action take $\lambda=0$. At an interior minimizer both action-bound multipliers vanish. Envelope subgradients thus have coordinates bounded by
\[
 \frac1d\left\{\frac{15}{2}+\beta\left(\frac{55}{8}+\frac{15}{32}\right)\right\}<\frac{15}{d}.
\]
Where the risk function is nonsmooth, the same bound holds for convex combinations of active risk gradients and valid KKT multipliers. Integrating bounded subgradients along a segment, with the one-sided limit at the boundary of $X$, proves \eqref{eq:lipschitz95}.
\end{proof}

\subsection{Exact one-date risk and outward construction}
Put $m=F(x,a,0)$, $R=1/32$ and $D=10d^{-1}\sum_jm_js_j$. Since $\sum_js_j=0$, the hinge in $g(m+sz)$ is constant in $z$. The target and cyclic terms give exactly
\begin{equation}\label{eq:terminalpoly95}
 g(m+sz)=g(m)+Dz+5z^2.
\end{equation}
Let $L=2R\alpha$ and
$h=\max\{-R,\min\{R-L,-D/10-L/2\}\}$.
The interval $[h,h+L]$ is a lowest-value interval of length $L$ for the convex quadratic. The maximizing AVaR density is $(1-\alpha)^{-1}$ on its complement and zero on that interval. The two tail moments are
\[
 M_1=-\frac{L(2h+L)}{2(2R-L)},\qquad
 M_2=\frac{2R^3-L(3h^2+3hL+L^2)}{3(2R-L)}.
\]
Hence the risk equals $g(m)+DM_1+5M_2$. This also follows from the threshold representation of AVaR, with the threshold chosen at the boundary of the retained lower interval. Endpoint minima and equal boundary values are included by the clipping rule.

For a lower witness, any dyadic $\widehat h\in[-R,R-L]$ defines an exactly feasible excluded interval and a feasible tail density. Differentiating its objective holds $\widehat h$ fixed. This distinction is essential: a derivative of a changing approximate tail need not be the derivative of a supporting objective. Upper values use an outward enclosure of the maximizing $h$, while lower witnesses use the separately retained feasible $\widehat h$. The lower witness therefore remains valid when the two differ.

At a fixed initial state, $F(x,u,z)$ is affine in $u$ and $z$. A convex continuation has a global affine minorant at every successor anchor. For a density constant on a bin, integration of this affine minorant equals its value at the midpoint. Its action dependence is affine. Adding the strong-convexity minorant of $h$ gives \eqref{eq:actionlower95}. A fitted continuation is used only where its nonnegative convex readout makes this argument valid; no analogous assertion is made for an unrestricted nonconvex fit.

For a value upper bound, convexity places a fitted continuation below the larger endpoint value on a bin. For $V_1$, the midpoint upper value plus $15R/n$ bounds the whole bin by \eqref{eq:lipschitz95}. If $l_i\leq v(z)\leq u_i$ on each bin, monotonicity of the original risk gives
\[
 \rho_\alpha\!\left(\sum_i l_i1_{B_i}\right)
 \leq\rho_\alpha(v)\leq
 \rho_\alpha\!\left(\sum_i u_i1_{B_i}\right).
\]
The step-function risks are finite upper-tail means because the bin probabilities are exactly equal and $n(1-\alpha)$ is an integer. Their computation does not assume that $v$ is constant on the bins.

\subsection{Complete policy and centered comparison}
\begin{proof}[Proof of Theorem~\ref{thm:continuous95}]
Every fixed density is feasible for the original risk supremum. Every continuation support is below $V_1$ at every successor. Accordingly its all-action quadratic minimum is below $V_0$, irrespective of which feasible proposal supplied the anchor. The binwise upper calculation dominates the value of the entire terminal controller; monotonicity and cash invariance charge at most $\beta\varepsilon_1$ for its terminal optimality allowance. It follows that
$\underline V_0\leq V_0\leq J_0^{\widehat\pi}\leq\overline J_0^{\widehat\pi}$.
Subtracting proves \eqref{eq:originalgap95}.

For the centered statement take any lower and upper continuation functions $\ell,u$ furnished by the preceding supports and bounds. Their differences from $f$ give admissible error functions. The exact centered identity cancels the common $\rho_\alpha(f)$ before any numerical enclosure. The preceding lower witness is a minorant of the resulting all-action expression; the preceding policy upper envelope is a majorant of the resulting policy expression. They are sufficient subsolution and supersolution inequalities in the general theorem. Computing them jointly is an implementation of that comparison, not a new objective. Loose global residual constants remain a valid alternative, but need not replace these sharper local witnesses.
\end{proof}

\paragraph{Whole-cell residual interpolation.}
On a state cell with vertices $v_i$, choose its multilinear convex weights $\omega_i(y)$. If $f$ and $V_1$ are convex, their vertex upper enclosures interpolate to upper bounds. A fixed affine support $L_V$ of $V_1$ satisfies
\[
 V_1(y)-f(y)\geq L_V(y)-\sum_i\omega_i(y)\overline f(v_i)
 =\sum_i\omega_i(y)\{L_V(v_i)-\overline f(v_i)\}.
\]
Its right side is at least the smallest vertex difference. Maximize that lower bound over the stored supports. Reversing $V_1$ and $f$, using upper values of $V_1$ and lower supports of $f$, gives the upper residual bound. The finite maximum or minimum over anchors and cells is taken after these whole-cell inequalities. No interpolation equality for the unknown function is asserted.

\paragraph{Strict tail exclusion.}
Suppose $u_i<\tau$ and the union $H$ of $k=n(1-\alpha)$ other bins has $v(z)\geq\tau$ throughout. It has probability at least $1-\alpha$. If a feasible maximizing density assigned positive mass to $B_i$, some capacity $q<(1-\alpha)^{-1}$ would remain on $H$. Moving equal mass from $B_i$ to $H$ strictly increases its risk objective, a contradiction. Equality yields only weak improvement and cannot support exclusion. Bin bounds used to construct $H$ must come from the unpruned calculation. The argument covers a continuum of densities without enumerating all of them.

